%%%%%%%%%%%%%%%%%%%%%%%%%%%%%%%%%%%%%%%%%%%%%%%%%%%%%%%%%%%%%%%%%%%%%%
% Lớp 10 | Đại số tổ hợp | Hoán vị. Chỉnh hợp. Tổ hợp
% Dạng 0D8?2-3: Bài toán đếm số tự nhiên
% Số câu: 68  (sinh tự động bởi tach_cau.py; từ nay thầy sửa trực tiếp file này)
%%%%%%%%%%%%%%%%%%%%%%%%%%%%%%%%%%%%%%%%%%%%%%%%%%%%%%%%%%%%%%%%%%%%%%

%%%%%%%%%%%%%%%% TRẮC NGHIỆM 4 PHƯƠNG ÁN %%%%%%%%%%%%%%%%

% nguồn: MaHoaNH_TracNghiem_Lop10
\begin{ex}%[0D8N2-3]%[ma:dae7c10e]
	 Từ tập hợp $E=\{1;2;3;4;5;6\}$ có thể tạo được bao nhiêu số tự nhiên có $4$ chữ số (các chữ số đôi một khác nhau)?
	 \choice
		 {\True $360$}
		 {$15$}
		 {$24$}
		{$4096$} 
	 \loigiai{
		Số  cách chọn số có $4$ chữ số khác nhau là $\mathrm A_{6}^4=360$.\\
	} 
\end{ex}

% nguồn: MaHoaNH_TracNghiem_Lop10
\begin{ex}%[0D8N2-3]%[ma:bdc3a3f1]
	Cho tập $X=\left\{1,2, 3,4, 5\right\}$. Viết được bao nhiêu số tự nhiên có $5$ chữ số khác nhau được lấy  từ tập $X$?
	 \choice
		 {$30!$}
		{$11!$} 
		 {\True $5!$}
		 {$6!$}
	\loigiai{ 
		Mỗi  số có $5$ chữ số khác nhau được lập từ $X$ là một hoán vị của $5$ phẩn tử của $X$. \\
		Vậy số các số có $5$ chữ số khác nhau được lấy từ tập $X$ là  $P_5=5!$.
	} 
\end{ex}

% nguồn: MaHoaNH_TracNghiem_Lop10
\begin{ex}%[0D8N2-3]%[ma:cd9c625e]
	Từ các chữ số $2$, $4$, $6$, $7$ người ta lập thành các số, mỗi số gồm $3$ chữ số. Số các số  lẻ lập được là
	 \choice
		 {\True $16$}
		{$24$} 
		{$6$} 
		{$27$} 
	 \loigiai{
		Có $1$ cách chọn chữ số lẻ là chữ số  $7$. Các chữ số còn lại có $4\cdot 4=16$ cách chọn.\\
		Vậy  có $1\cdot16=16$ số lẻ tạo thành.
	} 
\end{ex}

% nguồn: MaHoaNH_TracNghiem_Lop10
\begin{ex}%[0D8N2-3]%[ma:24490639]
	Có bao nhiêu số tự nhiên có  bốn chữ số?
	 \choice
		 {\True $9\,000$}
		 {$4\,536$}
		 {$5\,040$}
		{$10\,000$} 
	 \loigiai{
		Gọi số tự nhiên cần tìm là $n=\overline{a b c d}$, trong đó $a$, $b$, $c$, $d\in\left\{0,  1, 2, 3, 4, 5, 6, 7, 8, 9\right\}$ và $a\ne 0$.\\
		Ta có $a$ có $9$  cách chọn; $b$, $c$, $d$ mỗi số có $10$ cách chọn.\\
		 Vậy có cả thảy $9\cdot 10^3=9\,000$ số cần tìm.
	 }
\end{ex}

% nguồn: MaHoaNH_TracNghiem_Lop10
\begin{ex}%[0D8N2-3]%[ma:e132a1db]
	Cho các số $1$, $5$, $6$, $7$ có thể lập được bao nhiêu số  tự nhiên có $4$ chữ số và các chữ số khác nhau?
	\choice 
		 {$12$}
		{$64$} 
		 {\True $24$}
		{$256$} 
	\loigiai{ 
		Số cần tìm có dạng $\overline{abcd}$  với $a,b,c,d\in\{1,5,6,7\}$. Khi đó
		\begin{itemize}
			 \item $a$ có $4$ cách chọn.
			\item $b$ có $3$ cách chọn. 
			\item $c$  có $2$ cách chọn.
			\item $d$ có $1$  cách chọn.
		\end{itemize}
		Vậy có $4\cdot3\cdot2\cdot1=24$ số thỏa  mãn yêu cầu của đề bài.
	} 
\end{ex}

% nguồn: MaHoaNH_TracNghiem_Lop10
\begin{ex}%[0D8H2-3]%[ma:3da15b43]
	Từ $7$ chữ số $1$, $2$,  $3$, $4$, $5$, $6$, $7$ có thể lập được bao nhiêu số tự nhiên có $4$ chữ số đôi một khác nhau?
	\choice 
		{$\mathrm{C}_7^4$} 
		{\True $\mathrm{A}_7^4$} 
		 {$\mathrm{P}_7$}
		 {$7^4$}
	\loigiai{Từ $7$ chữ số $1$, $2$, $3$, $4$, $5$, $6$, $7$ có thể lập được  $\mathrm{A}_7^4$ số tự nhiên có $4$ chữ số đôi một khác nhau.}
\end{ex}

% nguồn: MaHoaNH_TracNghiem_Lop10
\begin{ex}%[0D8H2-3]%[ma:6be350b8]
	Có bao nhiêu số có ba chữ  số đôi một khác nhau mà các chữ số đó thuộc tập hợp $\left\{1;2;3;..;9\right\}$?
	 \choice
		{$\mathrm{C}_9^3$} 
		{$9^3$} 
		{\True $\mathrm{A}_9^3$} 
		{$3^9$} 
	 \loigiai{
		Số tự nhiên có ba chữ số đôi một khác nhau mà các chữ số đó thuộc tập hợp  $\left\{1;2;3;..;9\right\}$ là $\mathrm{A}_9^3$.
	 }
\end{ex}

% nguồn: MaHoaNH_TracNghiem_Lop10
\begin{ex}%[0D8H2-3]%[ma:10039a65]
	 Từ các chữ số $1$; $2$; $3$; $4$; $5$ có thể lập được bao nhiêu số tự nhiên có $5$ chữ số khác nhau?
	\choice 
		{\True $120$} 
		 {$5$}
		 {$625$}
		{$24$} 
	 \loigiai{
		 Có thể lập được $5!=120$ số tự nhiên thỏa mãn yêu cầu đề bài.
	} 
\end{ex}

% nguồn: MaHoaNH_TracNghiem_Lop10
\begin{ex}%[0D8H2-3]%[ma:d8b225a0]
	Từ $7$ chữ số $1$, $2$,  $3$, $4$, $5$, $6$, $7$ có thể lập được bao nhiêu số tự nhiên có $4$ chữ số khác nhau?
	\choice 
		 {$7!$}
		{$7^4$} 
		{\True  $\mathrm{A}_{7}^{4}$}
		 {$\mathrm{C}_{7}^{4}$}
	\loigiai{ 
		Mỗi số tự nhiên có $4$ chữ số khác nhau được lập là một cách chọn $4$ trong $7$ chữ số rồi sắp xếp  lại hay là một chỉnh hợp chập $4$ trong $7$ phần tử, do đó có thể lập được $\mathrm{A}_{7}^{4}$ số thỏa mãn yêu cầu đề bài.
	} 
\end{ex}

% nguồn: MaHoaNH_TracNghiem_Lop10
\begin{ex}%[0D8H2-3]%[ma:6194f0fc]
	Có bao  nhiêu số tự nhiên có $5$ chữ số, các chữ số khác $0$ và đôi một khác nhau?
	 \choice
		 {$5!$}
		{$9^5$} 
		 {$\mathrm{C}_9^5$}
		{\True $\mathrm{A}_9^5$} 
	 \loigiai{
		Mỗi cách lập số tự nhiên có $5$ chữ số, các chữ số khác $0$ và đôi một khác nhau là một chỉnh hợp chập  $5$ của $9$ phần tử.\\
		Vậy  có tất cả $\mathrm{A}_9^5$ số cần tìm.
	 }
\end{ex}

% nguồn: MaHoaNH_TracNghiem_Lop10
\begin{ex}%[0D8H2-3]%[ma:27507944]
	Từ các chữ số $0$; $1$; $2$; $3$; $4$; $5$; $6$; $7$; $8$ có thể lập được bao nhiêu số tự nhiên gồm có $5$ chữ số dạng $\overline{a_1 a_2 a_3 a_4 a_5}$ mà $a_1 < a_2 < a_3  < a_4 < a_5$.
	\choice 
		 {$52$}
		 {$26$}
		 {$282$}
		{\True $56$} 
	\loigiai{Vì $a_1 < a_2 < a_3 < a_4 < a_5$ nên $a_1$;  $a_2$; $a_3$; $a_4$; $a_5 \in \left\{1;2;\ldots;8\right\}$.\\
		Mỗi cách chọn $5$ số  thuộc tập trên cho ta duy nhất số thỏa yêu cầu bài toán.\\
		Vậy có $\mathrm{C}_8^5=56$ số thỏa mãn yêu cầu bài  toán.
	} 
\end{ex}

% nguồn: MaHoaNH_TracNghiem_Lop10
\begin{ex}%[0D8H2-3]%[ma:3e309714]
	Có bao nhiêu số tự nhiên có ba chữ số dạng $\overline{abc}$ với $a$, $b$,  $c$ $\in \left\{0;1; 2; 3; 4;5;6;7\right\}$ sao cho $a < b < c$?
	\choice 
		{\True  $35$}
		 {$56$}
		{$294$} 
		{$45$} 
	\loigiai{ 
		Vì số tự nhiên có ba chữ số dạng $\overline{abc}$ với $a$,  $b$, $c$ $\in \left\{0;1; 2; 3; 4;5;6;7\right\}$ sao cho $a < b < c$ ta có duy nhất một cách xếp chúng theo thứ tự tăng dần.\\
		Do đó các  số được lập là $\mathrm{C}_7^3=35$.
	} 
\end{ex}

% nguồn: MaHoaNH_TracNghiem_Lop10
\begin{ex}%[0D8H2-3]%[ma:8f83c30a]
	Cho $6$ chữ số $4$, $5$, $6$, $7$, $8$, $9$. Số các số tự nhiên  chẵn có $3$ chữ số khác nhau lập thành từ $6$ chữ số đó.
	\choice 
		 {$120$}
		{\True  $60$}
		 {$256$}
		{$216$} 
	 \loigiai{
		Gọi số cần tìm có dạng $\overline{abc}$. Chọn $c$ có 3 cách $(c\in \left\{4;  6; 8\right\})$.\\
		Chọn $\overline{ab}$  có $\mathrm{A}_5^2$ cách.\\
		Theo quy tắc nhân, có $3\cdot \mathrm{A}_5^2=60$ (số). 
	} 
\end{ex}

% nguồn: MaHoaNH_TracNghiem_Lop10
\begin{ex}%[0D8H2-3]%[ma:e7d1efb9]
	Có bao nhiêu số tự nhiên gồm $3$ chữ số đôi một khác nhau và khác $0$, biết rằng tổng của ba chữ số này bằng $8$? 
	\choice 
		{$24$} 
		{\True $12$} 
		{$6$} 
		{$18$} 
	 \loigiai{
		Gọi số cần  tìm là $\overline{a b c}$ với $a$, $b$, $c\in\{ 1;2;3;\ldots;9\}$ và đôi một khác nhau.\\
		 Do tổng của ba chữ số này bằng $8$ nên $a$, $b$, $c\in\{1; 2; 5\}$ hoặc $a$, $b$, $c\in\{ 1; 3; 4\}$.\\
		Do đó có $3!+3!=12$ số cần  tìm.
	 }
\end{ex}

% nguồn: MaHoaNH_TracNghiem_Lop10
\begin{ex}%[0D8H2-3]%[ma:7221b6a5]
	 Có bao nhiêu số tự nhiên có $5$ chữ số dạng $\overline{abcde}$ và thỏa mãn $a\le b<c\le d\le e$?
	\choice 
		 {$\mathrm{A}_9^5$}
		{$\mathrm{A}_{15}^5$} 
		{$\mathrm{C}_9^5$} 
		{\True $\mathrm{C}_{12}^5$} 
	\loigiai{ 
		Ta  có $1\le a<b+1<c+1<d+2<e+3\le 12$.\\
		Với mỗi bộ  số $a,b+1,c+1,d+2,e+3$ có $1$ số thỏa mãn đề bài.\\
		Vậy có $\mathrm{C}_{12}^5$ số thỏa mãn đề bài. 
	 }
\end{ex}

% nguồn: MaHoaNH_TracNghiem_Lop10
\begin{ex}%[0D8H2-3]%[ma:e6381bd4]
	Có bao nhiêu số tự nhiên gồm $6$ chữ số đôi một khác nhau, trong đó  có mặt chữ số $0$ nhưng không có mặt chữ số $1$?
	 \choice
		 {$151\,200$}
		 {$60\,480$}
		{$3\,024$} 
		 {\True $33\,600$}
	\loigiai{ 
		Gọi số cần lập là $\overline{abcdef}$ $(a\ne 0)$. \\ 
		Xếp chữ số $0$ vào các  vị trí từ $b$ đến $f$ có $5$ cách xếp. \\
		Còn lại $5$ vị trí, ta chọn $5$ trong $8$ chữ số để xếp  vào $5$ vị trí này có $\mathrm{A}_8^5$. \\
	Suy ra tất cả có  $5\cdot \mathrm{A}_8^5=33\,600$ số.}
\end{ex}

% nguồn: MaHoaNH_TracNghiem_Lop10
\begin{ex}%[0D8H2-3]%[ma:fd920bbe]
	Có bao nhiêu số gồm  $7$ chữ số khác nhau đôi một được lập bằng cách dùng $7$ chữ số $1$, $2$, $3$, $4$, $5$, $7$, $9$ sao cho hai chữ số chẵn không đứng liền nhau?
	 \choice
		 {\True $3\,600$}
		 {$1\,440$}
		 {$5\,040$}
		{$4\,320$} 
	\loigiai{ 
		Dùng $7$ chữ số đã cho, ta lập được tất cả $7!$ số có $7$ chữ số. \\ 
		Trong đó số các số có $2$ chữ số chẵn đứng liền nhau là $2\cdot  6!$ số. \\
		Vậy  số cần tìm có $7!-(2\cdot 6!)=3\,600$ số.
	 }
\end{ex}

% nguồn: MaHoaNH_TracNghiem_Lop10
\begin{ex}%[0D8H2-3]%[ma:f34fe591]
	Cho $5$ chữ số $1$, $2$, $3$, $4$, $6$.  Lập các số tự nhiên có $3$ chữ số đôi một khác nhau từ $5$ chữ số đã cho. Tính tổng của các số lập được.
	 \choice
		{$12\,321$} 
		{\True  $21\,312$}
		{$12\,312$} 
		 {$21\,321$}
	 \loigiai{
		Mỗi số số tự nhiên có $3$ chữ  số đôi một khác nhau từ $5$ chữ số $1$, $2$, $3$, $4$, $6$ là một chỉnh hợp chập $3$ của các chữ số này. Do đó, ta lập được $\mathrm{A}_{5}^{3}=60$ số.\\
		Do vai trò các số $1$, $2$, $3$, $4$, $6$ như nhau, nên số lần xuất hiện của mỗi chữ số trong các chữ số này ở mỗi hàng là như nhau và bằng $60:5=12$  lần.\\
		Vậy,  tổng các số lập được là
		$S=12\cdot( 1+2+3+4+6)\cdot( 100+10+1)=21\,312$. 
	} 
\end{ex}

% nguồn: MaHoaNH_TracNghiem_Lop10
\begin{ex}%[0D8H2-3]%[ma:74ddbe81]
	Có sáu quả cầu xanh đánh số từ $1$ đến $6$, năm quả cầu đỏ đánh số từ $1$ đến $5$ và bốn quả cầu vàng đánh số từ $1$ đến $4$. Hỏi có bao nhiêu  cách lấy ba quả cầu vừa khác màu vừa khác~số?
	\choice 
		{\True $64$} 
		 {$96$}
		{$128$} 
		{$32$} 
	\loigiai{ 
		Chọn $1$ quả cầu vàng có $4$  cách.\\
		Chọn $1$ quả cầu đỏ khác  số với quả cầu vàng vừa chọn có $4$ cách.\\
		Chọn $1$ quả cầu xanh khác số với quả cầu đỏ và quả cầu vàng vừa chọn  có $4$ cách.\\
		Vậy  có $4\cdot 4\cdot 4=64$ cách lấy thỏa mãn.
	} 
\end{ex}

% nguồn: MaHoaNH_TracNghiem_Lop10
\begin{ex}%[0D8H2-3]%[ma:3cbedf93]
	Có bao nhiêu số tự nhiên chẵn gồm $6$ chữ số khác  nhau trong đó chữ số đầu tiên là chữ số lẻ?
	 \choice
		 {$8\,400$}
		{\True  $42\,000$}
		{$60\,480$} 
		 {$33\,600$}
	\loigiai{ 
		Chọn vị trí thứ nhất là số  lẻ có $5$ cách.\\
		Chọn vị trí thứ sáu trong  các chữ số $\left\{0; 2; 4; 6; 8\right\}$ có $5$ cách. \\
		Chọn $4$ vị trí  còn lại có $\mathrm{A}_8^4$ cách. \\
		Vậy có $5\cdot 5\cdot \mathrm{A}_8^4=42000$ số. 
	} 
\end{ex}

% nguồn: MaHoaNH_TracNghiem_Lop10
\begin{ex}%[0D8H2-3]%[ma:24d2dc3e]
	Từ các chữ số  $1$, $2$, $3$, $5$, $7$ lập được bao nhiêu số tự nhiên có ba chữ số khác nhau và chia hết cho $2$?
	 \choice
		{\True $12$ số} 
		 {$20$ số}
		{$60$  số}
		 {$25$ số}
	\loigiai{ 
		Gọi số tự nhiên có ba chữ số có dạng  $\overline{abc}$ với $a$, $b$, $c\in \left\{1; 2; 3; 5; 7\right\}$. \\
		Vì $\overline{abc}$ chia hết cho $2$ nên có một cách  chọn $c=2$.\\
		Vì $a\ne b\ne  c$ nên số cách chọn $a$ và $b$ là $\mathrm{A}_4^2$. \\
		Mỗi cách chọn được  bộ ba số $a$, $b$, $c$ thỏa mãn điều kiện trên cho ta một số thỏa mãn yêu cầu bài toán. \\
		Vậy  có $1\cdot \mathrm{A}_4^2=12$ (số).
	} 
\end{ex}

% nguồn: MaHoaNH_TracNghiem_Lop10
\begin{ex}%[0D8H2-3]%[ma:e1065dd2]
	Có bao nhiêu số tự nhiên có $7$ chữ số khác nhau từng đôi một, trong đó  có $3$ chữ số $1$; $2$; $3$ và chữ số $2$ đứng liền giữa chữ số $1$ và chữ số $3$?
	\choice 
		{$2\,942$} 
		 {$5\,880$}
		 {$3\,204$}
		 {$7\,440$}
	\loigiai{ 
		Có $5$ cách chọn vị trí để xếp bộ  ba chữ số $\left\{1,2, 3\right\}$. \\
		Có $\mathrm{A}_7^4$ cách sắp xếp  $4$ chữ số được chọn từ tập hợp $\left\{0,4,5,6, 7,8, 9\right\}$. \\
		 Theo quy tắc nhân trường hợp này có $5\cdot 2!\mathrm{A}_7^4$ cách sắp xếp. \\
		 Trong các trường hợp ở trên có những trường hợp chữ số $0$ đứng đầu có $4\cdot 2!\cdot \mathrm{A}_6^3$ số dạng này. \\
		Vậy số các số tự nhiên thoả mãn bài ra là $5\cdot 2!\mathrm{A}_7^4-4\cdot  2!\mathrm{A}_6^3=7\,440$.
	} 
\end{ex}

% nguồn: MaHoaNH_TracNghiem_Lop10
\begin{ex}%[0D8H2-3]%[ma:7bbd286b]
	 Cho tập $E=\left\{0;1;2;3;4;5;6\right\}$. Hỏi có thể lập được bao nhiêu số tự nhiên gồm $4$ chữ số khác nhau chọn từ tập $E$ sao cho mỗi số chia hết cho $5$?
	\choice 
		 {$240$}
		{$1200$} 
		 {$100$}
		{\True  $220$}
	 \loigiai{
		Gọi số tự nhiên $4$ chữ số khác nhau  là $\overline{a b c d}$.
		\begin{enumerate}[TH1:]
			\item $d=5$, chọn $a$, $b$, $c$ lần lượt có $5$, $5$, $4$ cách.  Do đó $5\cdot5\cdot4=100$ số.
			\item $d=0$, chọn $a$, $b$, $c$ lần lượt có $6$, $5$, $4$ cách.  Do đó $6\cdot5\cdot4=120$ số.
		\end{enumerate}
		Vậy, theo quy tắc cộng có  $100+120=220$ số.
	 }
\end{ex}

% nguồn: MaHoaNH_TracNghiem_Lop10
\begin{ex}%[0D8H2-3]%[ma:3a218ea6]
	Xét biển số xe là dãy gồm $2$ chữ cái đứng trước và 4 chữ số đứng sau. Các chữ cái được lấy từ $26$ chữ cái $A$, $B$, $\ldots$, $Z$. Các chữ số được lấy từ $10$ chữ số $0$, $1$, $\ldots$, $9$. Có bao nhiêu biển số xe có hai chữ cái khác nhau đồng thời có đúng  hai chữ số lẻ và hai chữ số lẻ đó giống nhau.
	 \choice
		 {$81\,250$}
		 {$65\,000$}
		 {$13\,000$}
		 {$975\,000$}
	 \loigiai{
		Chọn chữ  cái thứ nhất có $26$ cách. \\
		Chọn chữ cái thứ  hai có $25$ cách.\\
		Chữ số chẵn thứ nhất chọn  từ $\{0,2, 4,6, 8\}$. Nên có $5$ cách.\\
		 Chữ số chẵn thứ hai chọn từ $\{0,2, 4,6, 8\}$. Nên có $5$ cách. \\
		Xếp $2$ chữ số  chẵn này vào $4$ vị trí, nên có $\mathrm{A}_4^2$ cách xếp. \\
		Hai chữ số lẻ giống  nhau có thể là $(1,1), (3,3), (5,5), (7,7), (9,9)$ và được xếp vào hai chỗ còn lại nên có $5$ cách. \\
		Vậy số biển xe cần tìm là  $\mathrm{A}_4^2\cdot 650\cdot 5^3=975\,000$ (biển số xe).
	 }
\end{ex}

% nguồn: MaHoaNH_TracNghiem_Lop10
\begin{ex}%[0D8H2-3]%[ma:221da868]
	Cho tập hợp $A=\left\{0; 1; 2; 3; 4; 5; 6; 7\right\}$. Hỏi từ tập $A$ có thể lập được bao nhiêu số tự nhiên gồm $5$ chữ số đôi một khác nhau sao  cho một trong $3$ chữ số đầu tiên phải bằng $1$.
	\choice 
		 {$2\,802$}
		{\True  $2\,280$}
		 {$65$}
		 {$2\,520$}
	\loigiai{ 
		Gọi số cần lập là $\overline{abcde}$ với $a$, $b$, $c$, $d$, $e\in A$ và $a>0$, các chữ số khác nhau. 
		\begin{itemize}
			\item TH1: Xét $a=1$. Số cách các chữ số còn lại là  $\mathrm{A}_7^4=840$.
			 \item TH2: Xét $a\ne 1$. \\
			Để chọn vị trí  cho chữ số $1$ có $2$ cách. \\
			Để hữ số $a$  có $6$ cách. \\
			Để các chữ số còn lại có  $\mathrm{A}_6^3$. \\
			Do đó có $2\cdot 6\cdot \mathrm{A}_6^3$ số lập được. 
		\end{itemize}
		Vậy có $\mathrm{A}_7^4+2\cdot 6\cdot \mathrm{A}_6^3=2\,280$ số thỏa mãn đề bài. 
	} 
\end{ex}

% nguồn: MaHoaNH_TracNghiem_Lop10
\begin{ex}%[0D8H2-3]%[ma:3c98dff7]
	Cho tập hợp $M=\left\{0; 1; 2; 3; 4; 5; 6; 7; 8\right\}$. Hỏi có thể lập được bao nhiêu số tự nhiên có ba chữ số đôi một khác nhau và chia hết cho $3$ mà  các chữ số thuộc tập M?
	\choice 
		{$180$} 
		 {$200$}
		{\True  $160$}
		{$140$} 
	\loigiai{ 
		Đặt $A=\left\{0,3, 6\right\}$, $B=\left\{1,4, 7\right\}$, $C=\left\{2,5,  8\right\}$. \\
		Gọi $\overline{abc}$ là số cần tìm ($a$, $b$, $c$  đôi một khác nhau và $\overline{abc}$ chia hết cho $3$). \\
		Giả sử a có thể bằng $0$. \\ 
		Trường hợp 1:  ba số $a, b, c$ cùng thuộc một tập $A$, $B$, $C$. Ta có $3. 3!$ số. \\
		Trường hợp  2: ba số $a, b, c$ mỗi số thuộc một tập $A$, $B$, $C$. Ta có $3!\cdot 3^3$ số. \\
		Vậy có $3\cdot 3!  + 3!\cdot 3^3= 180$ số. \\
		 Xét $a=0$. \\
		Ta có $\overline{0bc}$ chia hết cho $3$ khi hai số $b$, $c$ một số lấy ở $B$ một số lấy ở $C$ hoặc cả hai số $b$, $c$ khác $0$ và đều  thuộc tập $A$. \\
		Số cách lập là $3\cdot 3\cdot 2  +2 = 20$ số.\\
		Vậy số các số cần  tìm là $180-20 = 160$ (số).
	} 
\end{ex}

% nguồn: MaHoaNH_TracNghiem_Lop10
\begin{ex}%[0D8H2-3]%[ma:d8d2e398]
	Có bao nhiêu số tự nhiên  chẵn gồm $3$ chữ số khác nhau?
	 \choice
		{$500$} 
		 {$405$}
		 {$360$}
		{$328$} 
	\loigiai{ 
		Xét hai  trường hợp
		\begin{itemize}
			\item  TH1: Chữ số tận cùng là $0$ có $1$ cách chọn chữ số tận cùng. \\
			Có $\mathrm{A}_9^2$ cách chọn hai chữ số đầu.  \\
			Do  đó có $1\cdot \mathrm{A}_9^2 = 72$ số.
			\item TH2: Chữ số tận cùng là $2$, $4$,  $6$, $8$ có $4$ cách chọn chữ số tận cùng. \\
			Có $8$ cách chọn  chữ số đầu tiên. \\
			Có $8$ cách chọn  chữ số ở giữa. \\
			Do đó  có $4\cdot 8\cdot 8 = 256$ số.
		\end{itemize}
		Vậy có $72 + 256 = 328$ số thỏa mãn bài  toán.
	 }
\end{ex}

% nguồn: MaHoaNH_TracNghiem_Lop10
\begin{ex}%[0D8H2-3]%[ma:50c78638]
	Có bao nhiêu số tự nhiên  có $3$ chữ số đôi một phân biệt và chia hết cho $5$?
	\choice 
		 {$256$}
		 {$1458$}
		{$128$} 
		{\True  $136$}
	 \loigiai{
		Gọi $x=\overline{a b c}$ (với $a\ne b,\,\,b\ne c,\,\,c\ne a$) là số tự nhiên có  $3$ chữ số đôi một phân biệt và chia hết cho $5$. Vì $x \,\vdots\, 5$ nên $c\in\left\{0;5\right\}$.
		\begin{enumerate}[TH1:]
			 \item $c=0$.\\
			Chọn $c$ có $1$ cách.\\ 
			Chọn $a$ có $9$  cách ($a\ne 0$).\\
			Chọn $b$ có $8$ cách ($b\ne  0, b\ne a$).\\
			$\Rightarrow $ có $1\cdot9\cdot8=72$  số.
			\item $c=5$.\\ 
			Chọn $c$  có $1$ cách.\\
			Chọn $a$ có $8$ cách ($a\ne 5, a\ne 0$).\\ 
			Chọn $b$ có $8$ cách ($b\ne  5, b\ne a$).\\
			$\Rightarrow $ có $1\cdot8\cdot8=64$ số. 
		\end{enumerate}
		Theo quy tắc cộng, ta có tất cả $72+64=136$ số thỏa mãn yêu cầu  bài toán.
	} 
\end{ex}

% nguồn: MaHoaNH_TracNghiem_Lop10
\begin{ex}%[0D8H2-3]%[ma:4b50573c]
	Từ các chữ số $0$, $1$, $2$, $3$, $4$, $5$ có thể lập  được bao nhiêu số chẵn có $4$ chữ số đôi một khác nhau?
	 \choice
		{\True  $156$}
		{$240$} 
		 {$160$}
		 {$752$}
	 \loigiai{
		Gọi số cần tìm là $\overline{a b c d}$, trong đó $a\neq 0$ và $a$, $b$, $c$, $d$ đôi  một khác nhau.\\
		Vì $\overline{a b c d}$ là số chẵn  nên $d\in\left\{0,\text2,\text4\right\}$.
		\begin{enumerate}[TH1:]
			 \item Nếu $d=0$
			\begin{itemize}
				\item $d$ có $1$ cách chọn. 
				 \item $a$ có $5$ cách chọn.
				 \item $b$ có $4$ cách chọn.
				\item $c$ có  $3$ cách chọn.
			\end{itemize}
			Suy ra có $1\cdot5\cdot4\cdot3=60$ cách  chọn.
			 \item Nếu $d\ne 0$
			\begin{itemize}
				\item $d$ có $2$ cách chọn. 
				\item $a$ có $4$  cách chọn.
				\item $b$  có $4$ cách chọn.
				\item $c$  có $3$ cách chọn.
			\end{itemize}
			Suy ra có $2\cdot4\cdot4\cdot3=96$ cách  chọn.
		\end{enumerate}
		Vậy có $60+96=156$ số chẵn  có $4$ chữ số khác nhau.
	 }
\end{ex}

% nguồn: MaHoaNH_TracNghiem_Lop10
\begin{ex}%[0D8H2-3]%[ma:9b7e29d7]
	Có bao nhiêu số tự nhiên chia hết cho $11$, có $7$  chữ số đôi một khác nhau được thành lập từ các chữ số trong tập hợp $A=\left\{0; 1; 2; 3; 4; 5; 6\right\}$?
	 \choice
		 {$144$}
		{\True $288$} 
		 {$720$}
		{$4\,320$} 
	\loigiai{ 
		Gọi  số cần lập là $m=\overline{abcdefg}$.\\
		Vì $m$ chia hết cho $11$ nên $\left|a+c+e+g-(b+d+f)\right|$ là số chia hết cho $11$.  \\
		Ta có $a+b+c+d+e+f+g=0+1+\ldots +6=21$. \\ 
		Đặt  $\heva{
		& a+c+e+g=x  \\
		&  b+d+f=y} (x; y\in \mathbb{N})$\\
		 Ta có hệ $\hoac{
		& \heva{ & x+y=21 \\ 
		 & x-y=0 } \text{ (vô nghiệm)} \\
		 & \heva{
		&  x+y=21 \\
		& x-y=\pm 11.}}  \Leftrightarrow \hoac{
		 & \heva{
		 & x=16 \\
		 & y=5 } \\
		&  \heva{
		& x=5 \\ 
		&  y=16.}} $\\
		Nếu  $\heva{
		&  a+c+e+g=5 \\
		& b+d+f=16  } \Rightarrow $ Không tồn tại vì các chữ số $a$, $b$ đôi một khác nhau. \\
		 Nếu $\heva{
		 & a+c+e+g=16 \\
		& b+d+f=5} \Rightarrow \left\{b; d; f\right\}$  là các nhóm số $\left\{0; 1; 4\right\}$ hoặc $\left\{0; 2; 3\right\}$. \\
		Với mỗi trường hợp trên sẽ lập được $4!\cdot 3!=144$ số.  \\
		Vậy có $144\cdot  2=288$ số cần tìm.
	} 
\end{ex}

% nguồn: MaHoaNH_TracNghiem_Lop10
\begin{ex}%[0D8V2-3]%[ma:80a29158]
	Có tất cả bao nhiêu số tự nhiên có ba chữ số $\overline{a b c}$  sao cho $a$, $b$, $c$ là độ dài ba cạnh của một tam giác cân.
	 \choice
		{$45$} 
		{\True $165$} 
		{$81$} 
		{$216$} 
	 \loigiai{
		Không mất tính tổng quát, giả sử $a=b$ $\Rightarrow c<2a$ (Bất đẳng thức tam  giác).
		\begin{enumerate}[Th1:]
			\item  $a=b=1 \Rightarrow c<2\Rightarrow c=1$.
			\item $a=b=2 \Rightarrow c<4\Rightarrow c\in\{1;2;3\}$. 
			\item $a=b=3  \Rightarrow c<6\Rightarrow c\in\{1;2;3;4;5\}$.
			\item $a=b=4  \Rightarrow c<8\Rightarrow c\in\{1;2;3;4;5;\ldots;7\}$.
			\item $a=b\in\{ 5;6;7;8;9\} \Rightarrow  c\in\{ 1;2;3;\ldots;9\}$.
		\end{enumerate}
		Do đó có tất cả $1+3+5+7+9\cdot5=61$ bộ số thỏa mãn bài toán trong đó có $9$ bộ số là độ dài ba cạnh của một  tam giác đều và $52$ bộ số là độ dài của ba cạnh tam giác cân không đều.\\
		Với mỗi bộ ba số $a$, $b$, $c$ là độ dài ba cạnh của tam giác cân, ta có $3$ cách sắp  xếp để tạo thành một số có ba chữ số.\\
		 Vậy số các số cần tìm là $9+52\cdot3=165$ số.
	 }
\end{ex}

% nguồn: MaHoaNH_TracNghiem_Lop10
\begin{ex}%[0D8V2-3]%[ma:79326b9e]
	Có bao nhiêu số tự nhiên lẻ có $6$ chữ số và  chia hết cho $9$?
	\choice 
		 {$30\,000$}
		 {$40\,000$}
		{\True  $50\,000$}
		{$60\,000$} 
	\loigiai{ 
		 \textbf{Cách 1:}\\
		Gọi số có $6$ chữ số và chia hết cho  $9$ là $\overline{a b c d e f}$.\\
		Số tự nhiên lẻ nên $f$  có số cách chọn là $5$ cách.\\
		Số cách chọn $b$, $c$, $d$, $e$ là  $10^4$ cách.\\
		Gọi $r$ là số dư khi tổng $\left(b+c+d+e+f\right)$  chia cho $9$.\\
		Để $\overline{a b c d e f}$ chia  hết cho $9$ thì $r+a\,\vdots\, 9$. Mà $0<r+a<18 \Rightarrow r+a=9 \Rightarrow a=9-r \Rightarrow a$ có $1$ cách chọn.\\
		Vậy số các số tự nhiên lẻ có $6$ chữ số  và chia hết cho $9$ là $5\cdot10^4\cdot1=50\,000$ số.\\
		 \textbf{Cách 2:}\\
		Số lẻ nhỏ nhất có $6$ chữ  số và chia hết cho $9$ là $100\,017$.\\
		Số lẻ lớn nhất có $6$ chữ số và chia hết cho  $9$ là $999\,999$.\\
		Hai số liên tiếp là số lẻ có $6$  chữ số và chia hết cho $9$, cách nhau $18$ đơn vị.\\
		Vậy số các số tự nhiên lẻ  có $6$ chữ số và chia hết cho $9$ là $\dfrac{999\,999-100\,017}{18}+1=50\,000$ số.
	} 
\end{ex}

% nguồn: MaHoaNH_TracNghiem_Lop10
\begin{ex}%[0D8V2-3]%[ma:5d955be2]
	Từ $2$ chữ số $1$ và $8$ lập được bao  nhiêu số tự nhiên có $8$ chữ số sao cho không có $2$ chữ số $1$ đứng cạnh nhau?
	\choice 
		{$54$} 
		{$110$} 
		 {\True $55$}
		{$108$} 
	 \loigiai{
		Trường hợp 1: Có $8$ chữ số $8$  nên có $1$ số.\\
		Trường hợp 2: Có $1$ chữ số $1$, $7$ chữ số $8$ nên  có $8$ cách xếp chữ số $1$ nên có $8$ số.\\
		Trường hợp 3: Có $2$ chữ số $1$,  $6$ chữ số $8$.\\
		Xếp $6$ số  $8$ ta có $1$ cách.\\
		Từ $6$ số  $8$ ta có có $7$ chỗ trống để xếp $2$ số $1$.\\
		Nên  ta có $\mathrm{C}_7^2=21$ số.\\
		Trường hợp 4: Có $3$ chữ số $1$, $5$ chữ số  $8$.\\
		Tương tự TH3, từ $5$ chữ số $8$ ta có $6$ chỗ trống để xếp $3$ chữ số  $1$.\\
		Nên  có $\mathrm{C}_6^3=20$ số.\\
		Trường hợp 5: Có  $4$ chữ số $1$, 4 chữ số $8$.\\
		Từ $4$ chữ số $8$ ta có $5$ chỗ trống để xếp $4$ chữ số $1$.\\ 
		 Nên có $\mathrm{C}_5^4=5$.\\
		Vậy có $1+8+21+20+5=55$  số.
	} 
\end{ex}


%%%%%%%%%%%%%%%% ĐÚNG/SAI %%%%%%%%%%%%%%%%

% nguồn: MaHoaNH_DungSai_Lop10
\begin{ex}%[0D8H2-3]%[ma:af2308ab]
	Chọn khẳng định đúng trong các khẳng định sau 
	\choiceTF 
		{\True Có $15120$ số tự nhiên có $5$ chữ số, các chữ số khác $0$  và đôi một khác nhau}
		{Có $30240$ số tự nhiên có $5$ chữ số, các chữ  số khác $0$ và đôi một khác nhau}
		{\True Có $3024$ số tự nhiên có $4$  chữ số, các chữ số khác $0$ và đôi một khác nhau}
		{Có $5040$ số tự nhiên có $4$ chữ số, các chữ  số khác $0$ và đôi một khác nhau}
	\loigiai{ 
		\begin{itemize}
			\item Mỗi số tự nhiên tự nhiên có $5$ chữ số khác nhau được lập từ $9$ chữ số $1, 2, 3, 4, 5, 6, 7, 8, 9$ là một chỉnh hợp chập $5$ của $9$ phần tử. Vậy có $\mathrm{A}_5^9=15120$ số tự nhiên có  $5$ chữ số , các chữ số khác $0$ và đôi một khác nhau.
			\item Mỗi số  tự nhiên tự nhiên có $4$ chữ số khác nhau được lập từ $9$ chữ số $1, 2, 3, 4, 5, 6, 7, 8, 9$ là một chỉnh hợp chập $4$ của $9$ phần tử. Vậy có $\mathrm{A}_4^9=3024$ số tự nhiên có $4$ chữ số, các chữ số khác $0$ và đôi một khác nhau.
		\end{itemize}
	 }
\end{ex}

% nguồn: MaHoaNH_DungSai_Lop10
\begin{ex}%[0D8H2-3]%[ma:d2d0e6d8]
	Từ các chữ số thuộc tập hợp $X=\{1;2;3;4;5\}$, người ta  lập các số tự nhiên có bốn chữ số.
	 \choiceTF
		{\True Có $\mathrm{A}_5^3$ số có các chữ số đôi một khác  nhau}
		{\True Có  $4\mathrm{A}_4^3$ số có các chữ số đôi một khác nhau và có mặt chữ số $1$}
		{Có $3\mathrm{C}_5^3$ số có các  chữ số đôi một khác nhau và là số lẻ}
		{Có $4\mathrm{A}_5^3$ số có các chữ số đôi một khác nhau và là số chia  hết cho $5$}
	 \loigiai{
		\begin{itemchoice}
			\itemch Đúng. Vì các chữ số đôi một khác nhau là chỉnh hợp chập $3$ của  $5$, nên có $\mathrm{A}_5^3$ cách.
			\itemch Đúng. Vì có $4$ vị trí số $1$, ba vị trí còn lại là chỉnh hợp chập  $3$ của $4$, nên ta có $4\mathrm{A}_4^3$ cách.
			\itemch  Sai. Vì ta có $3$ cách chọn số tận cùng là số lẻ, nên $3$ vị trí còn lại là chỉnh hợp chập $3$ của $4$, nên có $3\mathrm{A}_4^3$ cách.
			 \itemch Sai. Vì ta có $1$ cách chọn số tận cùng là số $5$, nên $3$ vị trí còn lại là chỉnh hợp chập $3$ của $4$, nên có $\mathrm{A}_4^3$ cách.
		\end{itemchoice}
	} 
\end{ex}

% nguồn: MaHoaNH_DungSai_Lop10
\begin{ex}%[0D8H2-3]]%[ma:4b7b8543]
	Chọn khẳng định đúng trong  các khẳng định sau
	\choiceTF 
		{Từ  các chữ số $1$, $2$, $3$, $4$ có thể lập được $12$ số tự nhiên có $4$ chữ số}
		{\True Từ các chữ số $1$, $2$, $3$, $4$  có thể lập được $24$ số tự nhiên có $4$ chữ số đôi một khác nhau}
		 {\True Từ các chữ số $1$, $2$, $3$, $4$ có thể lập được $12$ số tự nhiên có $2$ chữ số đôi một khác nhau}
		{Từ các chữ số $1$, $2$, $3$, $4$ có thể lập được $24$ số tự nhiên có $4$ chữ  số giống nhau}
	\loigiai{ 
		\begin{itemchoice}
			\itemch Từ các chữ số $1$, $2$, $3$, $4$ có thể lập  được $4\cdot 4=16$ số tự nhiên có $2$ chữ số.
			\itemch Mỗi số tự nhiên tự nhiên có $4$ chữ số khác nhau được  lập từ $4$ chữ số đã cho là một hoán vị của $4$ phần tử. Vậy có $4!=24$ số tự nhiên có $4$ chữ số đôi một khác nhau.
			\itemch Mỗi số tự nhiên tự nhiên có $2$ chữ  số khác nhau được lập từ $4$ chữ số đã cho là một chỉnh hợp chập $2$ của $4$ phần tử. Vậy có $\mathrm{A}_2^4=12$ số tự nhiên có $2$ chữ số đôi một khác nhau.
			\itemch Từ các chữ số $1$, $2$, $3$, $4$ có thể lập  được $4$ số tự nhiên có $4$ chữ số giống nhau là $1111$, $2222$, $3333$ và $4444$.
		\end{itemchoice}
	 }
\end{ex}

% nguồn: MaHoaNH_DungSai_Lop10
\begin{ex}%[0D8H2-3]%[ma:37f40e40]
	Từ các chữ số $A=\left\{1;3;4;7;8\right\}$ lập thành số có $3$  chữ số đôi một khác nhau.
	 \choiceTF
		{\True Có $24$  số mà trong đó không có mặt chữ số $1$}
		{\True Có $36$  số là số lẻ}
		{Có $36$ số là  số chẵn}
		{Trong các số được tạo thành thì số chẵn nhiều hơn số lẻ} 
	\loigiai{ 
		\begin{itemchoice}
			\itemch 
			Do số  được tạo thành không có mặt chữ số $1$ nên ta có tập hợp $B=\left\{3;4;7;8\right\}$.\\
			Ta chọn $3$ chữ số bất kì trong tập hợp $B$ để có được số các số có $3$ chữ số đôi một khác nhau không có  mặt chữ số $1$. Số các số cần tìm là $\mathrm{A}_4^3=24$ (số).
			\itemch 
			Gọi $\overline{abc}$ là số có  $3$ chọn $c\in \left\{1;3;7\right\}$ có $3$ cách.\\
			Chọn  $\overline{ab}$ có $\mathrm{A}_4^2$ cách.\\
			Áp dụng quy tắc  nhân ta có $3\cdot \mathrm{A}_4^2=36$ số cần tìm.
			\itemch 
			Gọi  $\overline{abc}$ là số có $3$ chữ số cần tìm.\\
			Do  số cần tìm là số chẵn nên ta chọn $c\in \left\{4;8\right\}$ có $2$ cách.\\
			Chọn $\overline{ab}$ có $\mathrm{A}_4^2$  cách.\\
			Áp dụng quy tắc nhân ta có $2\cdot \mathrm{A}_4^2=24$  số cần tìm.
			 \itemch
		\end{itemchoice}
	 }
\end{ex}

% nguồn: MaHoaNH_DungSai_Lop10
\begin{ex}%[0D8H2-3]%[ma:200c5841]
	Cho các  chữ số $0$; $2$; $3$; $8$; $9$.
	 \choiceTF
		{\True Số các số tự nhiên có  $2$ chữ số khác nhau được lập từ năm chữ số trên là $16$}
		{Số  các số tự nhiên có $2$ chữ số khác nhau được lập từ năm chữ số trên là $20$}
		{Số các số tự nhiên lẻ có $2$ chữ số khác nhau được lập từ năm  chữ số trên là $10$}
		{Số các số tự nhiên lẻ có $2$ chữ số khác nhau được lập từ năm chữ số trên là  $8$}
	 \loigiai{
		\begin{itemchoice}
			 \itemch
			Gọi số tự nhiên có hai chữ số  khác nhau được lập từ năm chữ số trên là $\overline{ab}$.\\
			Số cách  chọn số $a$ khác $0$ là $4$ cách.\\
			Số cách chọn số $b$ khác $a$ là  $4$ cách.\\
			Vậy có $4\cdot 4=16$ số thỏa  mãn yêu cầu bài toán.
			 \itemch
			Số các số tự nhiên có $2$ chữ số  khác nhau được lập từ năm chữ số trên là $16$.
			\itemch 
			Gọi  số tự nhiên lẻ có hai chữ số khác nhau được lập từ năm chữ số trên là $\overline{ab}$.\\
			Số  cách chọn số $b$ là $2$ cách.\\
			Số cách chọn số $a$ khác $0$ và khác $b$  là $3$ cách.\\
			Vậy có $2\cdot  3=6$ số thỏa mãn yêu cầu bài toán.
			\itemch 
			Số các số tự nhiên lẻ có $2$ chữ số  khác nhau được lập từ năm chữ số trên là $6$.
		\end{itemchoice}
	} 
\end{ex}

% nguồn: MaHoaNH_DungSai_Lop10
\begin{ex}%[0D8V2-3]%[ma:6c9f2eba]
	Cho tập $A=\{0;1;2;3;4;5\}$. Lập các số tự  nhiên từ tập $A$.
	 \choiceTF
		{Số các số có hai  chữ số là $25$}
		{Số các số lẻ  có hai chữ số khác nhau là $15$}
		{\True Số các số chẵn có ba chữ số khác nhau là $52$} 
		{Gọi $T$ là số  các số số tự nhiên gồm $5$ chữ số khác nhau, trong đó nhất thiết phải có chữ số $0$ và $3$. Tổng các chữ số của $T$ bằng $8$}
	 \loigiai{
		\begin{itemchoice}
			\itemch  Số các số có hai chữ số là $5\cdot 6=30$.
			\itemch Số các  số lẻ có hai chữ số là $3\cdot 4=12$.
			\itemch Số các số chẵn có ba chữ số khác nhau là $1\cdot \mathrm{A}_{5}^{2}+2\cdot  4\cdot 4=52$.
			\itemch Gọi số cần tìm là  $\overline{abcde}\,(a\neq 0)$.\\
			Tìm số các số có $5$ chữ số khác nhau mà có  mặt $0$ và $3$ không xét đến vị trí $a$.\\
			Xếp $0$ và $3$ vào  $5$ vị trí có $\mathrm{A}_{5}^{2}$ cách.\\
			$3$ vị trí  còn lại có $\mathrm{A}_{4}^{3}$ cách.\\
			Suy ra có $\mathrm{A}_{5}^{2}\cdot \mathrm{A}_{4}^{3}$  số.\\
			Tìm số các số có $5$  chữ số khác nhau mà có mặt $0$ và $3$ với $a=0$.\\
			Xếp $3$  có $4$ cách và $3$ vị trí còn lại có $\mathrm{A}_{4}^{3}$ cách.\\
			Suy ra có $4\cdot \mathrm{A}_{4}^{3}$  số.\\
			Vậy số các số  cần tìm thoả mãn yêu cầu bài toán là $\mathrm{A}_{5}^{2}\cdot \mathrm{A}_{4}^{3}-4\cdot \mathrm{A}_{4}^{3}=384$.
		\end{itemchoice}
	} 
\end{ex}

% nguồn: MaHoaNH_DungSai_Lop10
\begin{ex}%[0D8V2-3]%[ma:35ff2575]
	Từ  các chữ số $1$; $3$; $5$; $7$; $9$.
	 \choiceTF
		{\True Có thể lập được $120$ số có  năm chữ số khác nhau}
		{Có thể lập được  $90$ số có bốn chữ số khác nhau trong đó phải có mặt chữ số $5$}
		{\True Có thể lập được $48$ số có  năm chữ số khác nhau trong đó chữ số $7$; $9$ luôn đứng cạnh nhau}
		{\True Có  thể lập được $20$ số có ba chữ số khác nhau và chia hết cho $3$}
	\loigiai{ 
		\begin{itemchoice}
			\itemch 
			Mỗi số tự nhiên tạo thành là một hoán vị của năm số đã  cho. Vậy có $\mathrm{P}_5=5!=120$ (số).
			 \itemch
			Số cách lập số có  bốn chữ số khác nhau từ các chữ số $1$; $3$; $5$; $7$; $9$.\\
			Gọi số cần tìm có dạng $\overline{abcd}$ trong đó $a$, $b$,  $c$, $d$ khác nhau.\\
			Chọn $a$ có $5$ cách, chọn $b$ có $4$ cách, chọn $c$ có $3$ cách, chọn $d$  có $2$ cách do đó có $5\cdot 4\cdot 3\cdot 2=120$ (số).\\
			Số cách lập số có bốn chữ số khác nhau trong đó không có mặt chữ số $5$ là  $4!=24$ (số).\\
			Vậy có $120-24=96$ số có bốn chữ số khác nhau  trong đó phải có mặt chữ số $5$.
			 \itemch
			Xét $A=\left\{1;3;5;79\right\}$ và  $B=\left\{1,3,5,97\right\}$. \\
			Từ  mỗi tập $A$ và $B$ ta lập được $4!$ số có năm chữ số khác nhau trong đó số $7$ và $9$ luôn đứng cạnh nhau. \\
			Vậy có  $2\cdot 4!=48$ (số).
			\itemch 
			Các bộ ba số lấy từ  $5$ số đã cho có tổng chia hết cho $3$ là $\left\{1;3;5\right\},\left\{1;5;9\right\},\left\{3;5;7\right\}$,\\$\left\{5;7;9\right\}$.
			Từ mỗi bộ  này ta lập được $3!$ số có ba chữ số khác nhau và chia hết cho $3$.\\
			Vậy có tất cả $4\cdot  3!=24$ số.
		\end{itemchoice}
	 }
\end{ex}

% nguồn: MaHoaNH_DungSai_Lop10
\begin{ex}%[0D8V2-3]%[ma:6c916702]
	Cho tập  hợp $A=\left\{1;2;3;4;5;6;7\right\}$.
	 \choiceTF
		{Số các số tự  nhiên có bốn chữ số khác nhau được lập từ $A$ là $4^7$}
		{\True Số các số tự nhiên có ba chữ số lẻ khác nhau được lập từ  $A$ là $\mathrm{A}_4^3$}
		{Số các số tự nhiên chẵn  có bốn chữ số khác nhau được lập từ $A$ là $4\cdot\mathrm{A}_6^3$}
		{\True Số các số tự nhiên có bốn chữ số khác nhau trong đó có một chữ số chẵn và ba chữ số lẻ  được lập từ $A$ là $12\cdot\mathrm{A}_4^3$}
	 \loigiai{
		\begin{itemchoice}
			 \itemch
			Mỗi  số tự nhiên có bốn chữ số khác nhau được lập từ $A$ là một chỉnh hợp chập $4$ của $7$.\\
			Do đó số các số tự  nhiên có bốn chữ số khác nhau được lập từ $A$ là $\mathrm{A}_7^4=840$ số.
			\itemch 
			Để lập một  số tự nhiên có ba chữ số lẻ khác nhau được lập từ $A$ ta lấy ba chữ số lẻ từ bốn chữ số lẻ của $A$ và sắp xếp ba chữ số lấy được vào ba vị trí hàng đơn vị, hàng chục và hàng trăm. \\
			Mỗi kết quả của quá trình  trên là một chỉnh hợp chập $3$ của $4$ nên số các số tự nhiên có ba chữ số lẻ khác nhau được lập từ $A$ là $\mathrm{A}_4^3=24$ số.
			\itemch 
			Để lập một số tự nhiên chẵn có bốn chữ số  khác nhau được lập từ $A$ ta thực hiện lần lượt $2$ bước sau
			\begin{enumerate}[\bf Bước 1.]
				\item Chọn chữ số hàng đơn vị là chữ số chẵn, có $3$ cách chọn. 
				\item Lấy ba trong sáu chữ số còn lại của $A$ và sắp xếp vào các vị  trí hàng chục, hàng trăm, hàng nghìn. \\
				Mỗi kết quả của quá trình này  là một chỉnh hợp chập $3$ của $6$ nên có $\mathrm{A}_6^3$ cách chọn.
			\end{enumerate}
			Áp dụng quy tắc nhân ta  được số các số tự nhiên chẵn có bốn chữ số khác nhau được lập từ $A$ là $3\cdot\mathrm{A}_6^3=360$ số.
			\itemch 
			Để lập một số tự nhiên có bốn chữ số khác nhau trong đó có một chữ số chẵn và ba chữ số lẻ được lập từ $A$ ta thực hiện lần lượt $2$ bước  sau
			\begin{enumerate}[\bf Bước 1.]
				\item Chọn một chữ số chẵn và xếp vào một trong bốn vị trí hàng đơn vị,  hàng chục, hàng trăm, hàng nghìn. Có $3\cdot 4=12$ cách.
				\item Lấy ba chữ số lẻ trong bốn chữ số lẻ và  xếp vào bốn vị trí còn lại. Có $\mathrm{A}_4^3$ cách.
			\end{enumerate}
			Áp dụng quy tắc nhân ta được số các số tự nhiên có bốn chữ số khác  nhau trong đó có một chữ số chẵn và ba chữ số lẻ được lập từ $A$ là $12\cdot\mathrm{A}_4^3=288$ số.
		\end{itemchoice}
	} 
\end{ex}

% nguồn: MaHoaNH_DungSai_Lop10
\begin{ex}%[0D8V2-3]%[ma:cf5abe80]
	Cho tập $S=\left\{0;1;2;3;4;5;6\right\}$. 
	\choiceTF 
		{\True Có $6\cdot 6!$ số tự nhiên có $7$ chữ số đôi một khác nhau được lấy  từ tập $S$}
		{Có $144$  số tự nhiên có $7$ chữ số đôi một khác nhau được lấy từ tập $S$ sao cho $3$ chữ số $1$; $2$; $3$ luôn đứng cạnh nhau}
		{\True  Có $6!$ số tự nhiên có $6$ chữ số khác nhau được lấy từ tập $S\setminus \{0\}$}
		{Có $3\cdot 5\cdot 5!$ số tự nhiên có  $7$ chữ số đôi một khác nhau sao cho số đó là số chẵn}
	 \loigiai{
		\begin{itemchoice}
			\itemch 
			Gọi số cần tìm có dạng $\overline{a_1a_2 a_3a_4a_5 a_6 a_7}$ trong đó $a_1$;  $a_2$; $a_3$; $a_4$; $a_5$; $a_6$; $a_7 \in S$.\\
			 Chọn $a_1 (a_1 \ne 0)$ có $6$ cách chọn.\\
			Ta có $a_2$; $a_3$; $a_4$; $a_5$; $a_6$; $a_7$ có số  cách chọn là số hoán vị của $6$ phần tử $6!$.\\
			Vậy có $6\cdot 6!$  số.
			\itemch 
			Gọi số cần tìm có dạng $\overline{a_1 a_2 a_3 a_4 a_5 a_6 a_7}$  trong đó $a_1$; $a_2$; $a_3$; $a_4$; $a_5$; $a_6$; $a_7 \in S$.
			\begin{enumerate}[\bf Trường hợp 1.]
				\item Chọn $a_1$; $a_2$; $a_3 \in \left\{1;2;3\right\}$, $(a_1 \ne 0)$ có  $3!$ cách chọn.\\
				Chọn $a_4$; $a_5$; $a_6$; $a_7$ có số cách chọn là số  hoán vị của $4$ phần tử còn lại $4!$ cách chọn.\\
				 Do vậy ta được $3!\cdot 4!=144$ số.
				\item Các số $1$; $2$; $3$ nằm ở  ba trong $4$ vị trí $a_4$; $a_5$; $a_6$; $a_7$ có $4\cdot 3\cdot 2=24$ cách sắp xếp.\\
				Chọn $a_1  \in \left\{4;5;6\right\}$ có $3$ cách chọn.\\
				Còn $3$ vị trí còn lại có số cách chọn  là số hoán vị của $3$ phần tử còn lại từ tập $S$ có $3!$ cách chọn.\\
				Do vậy  ta có $24\cdot 3\cdot 3!=432$ số.
			\end{enumerate}
			Tổng cộng có  $576$ số.
			\itemch 
			Gọi số cần tìm có dạng $\overline{a_1 a_2 a_3 a_4 a_5 a_6}$ trong đó $a_1$; $a_2$; $a_3$; $a_4$; $a_5$;  $a_6 \in S\setminus \{0\}$.\\
			Ta có $a_1$; $a_2$;  $a_3$; $a_4$; $a_5$; $a_6$ có số cách chọn là số hoán vị của $6$ phần tử $6!$.\\
			Do vậy ta có $6!$  số.
			 \itemch
			Gọi số cần tìm có dạng $\overline{a_1a_2a_3a_4a_5a_6  a_7}$ trong đó $a_1;a_2;a_3;a_4;a_5;a_6;a_7 \in S$.
			\begin{enumerate}[\bf Trường hợp 1.]
				\item Chọn $a_7=0$  có $1$ cách chọn.\\
				Chọn $a_1$; $a_2$;  $a_3$; $a_4$; $a_5$; $a_6$ có số cách chọn là số hoán vị của $6$ phần tử $6!$.\\
				 Do vậy ta có $6!$ số.
				\item Chọn $a_7 \in \left\{2,4,6\right\}$ có $3$  cách chọn.\\
				Chọn $a_1 (a_1 \ne 0; a_1  \ne a_7)$ có $5$ cách chọn.\\
				Chọn $a_2;a_3;a_4;a_5;a_6$ có số cách chọn là số  hoán vị của $5$ phần tử $5!$.\\
				Do vậy  ta có $3\cdot 5\cdot 5!$ số.
			\end{enumerate}
			 Vậy tổng có $6!+3\cdot 5\cdot 5!$.
		\end{itemchoice}
	 }
\end{ex}

% nguồn: MaHoaNH_DungSai_Lop10
\begin{ex}%[0D8V2-3]%[ma:1a5bab9f]
	Từ các chữ số $0$; $1$; $2$; $3$;  $4$; $5$; $6$; $7$; $8$; $9$.
	\choiceTF 
		{\True Có thể lập được  $4\,536$ số có bốn chữ số khác nhau}
		{\True Có thể lập được  $13\,440$ số lẻ có năm chữ số khác nhau}
		 {Có thể lập được $2\,240$ số chẵn có bốn chữ số khác nhau}
		{\True Có thể  lập được $36\,960$ số lẻ gồm sáu chữ số khác nhau và nhỏ hơn $600\,000$}
	\loigiai{ 
		\begin{itemchoice}
			\itemch 
			Gọi số  cần tìm có dạng $\overline{abcd}$, trong đó $a$, $b$, $c$, $d$ khác nhau.\\
			$a\ne 0\Rightarrow a$  có $9$ cách chọn.\\
			Chọn $\overline{bcd}$ có  $\mathrm{A}_9^3$ cách.\\
			Theo quy tắc nhân,  ta có $4\,536$ số.
			\itemch 
			Gọi số cần tìm có dạng  $\overline{abcde}$, trong đó $a$, $b$, $c$, $d$, $e$ khác nhau.\\
			Vì số cần tìm  là số lẻ nên $e\in \left\{1;3;5;7;9\right\}$. Do đó, $e$ có $5$ cách chọn.\\
			$\heva{&a\ne 0\\&a\ne e} \Rightarrow a$ có $8$  cách chọn.\\
			Chọn $\overline{bcd}$ có $\mathrm{A}_8^3$  cách.\\
			Theo quy tắc nhân, ta có $13\,440$ số. 
			\itemch 
			Gọi số cần tìm có dạng $\overline{abcd}$, trong đó $a$,  $b$, $c$, $d$ khác nhau.\\
			Vì  số cần tìm là số chẵn nên $d\in \left\{0;2;4;6;8\right\}$.
			\begin{enumerate}[\bf Trường hợp 1.]
				\item  $d=0$ suy ra $d$ có $1$ cách chọn.\\
				$\heva{&a\ne 0\\&a\ne d\left(d=0\right)} \Rightarrow a$ có $9$ cách chọn.\\ 
				Chọn  $\overline{bc}$ có $\mathrm{A}_8^2$ cách.\\
				Theo quy tắc nhân, ta  có $504$ số.
				\item $d\in \left\{2;4;6;8\right\}$  suy ra $d$ có $4$ cách chọn.\\
				$\heva{&a\ne 0\\&a\ne d; \; \left(d\in \left\{2;4;6;8\right\}\right)} \Rightarrow a$ có  $8$ cách chọn.\\
				Chọn $\overline{bc}$  có $\mathrm{A}_8^2$ cách.\\
				Theo quy tắc nhân, ta  có $1\,792$ số.
			\end{enumerate}
			Theo  quy tắc cộng, ta có $2\,296$ số.
			 \itemch
			Gọi số cần tìm có dạng $\overline{abcdef}$, trong đó $a$, $b$, $c$, $d$, $e$, $f$  khác nhau.\\
			Vì số cần  tìm là số lẻ và nhỏ hơn $600\,000$ nên $\heva{&a\in \left\{1;2;3;4;5\right\}\\&f\in \left\{1;3;5;7;9\right\}}$.
			\begin{enumerate}[\bf Trường hợp 1.]
				\item $a\in \left\{1;3;5\right\}$, $a$ có $3$  cách chọn.\\
				$\heva{&f\in \left\{1;3;5;7;9\right\}\\&f\ne a} \Rightarrow f$  có $4$ cách chọn.\\
				 Chọn $\overline{bcde}$ có $\mathrm{A}_8^4$ cách.\\
				Theo quy tắc  nhân, ta có $20\,160$ số.
				\item $a\in \left\{2;4\right\}$, $a$ có $2$ cách  chọn.\\
				$\heva{&f\in \left\{1;3;5;7;9\right\}\\&f\ne  a} \Rightarrow f$ có $5$ cách chọn.\\
				Chọn $\overline{bcde}$  có $\mathrm{A}_8^4$ cách.\\
				Theo quy tắc nhân,  ta có $16\,800$ số.
			\end{enumerate}
			Theo quy tắc cộng, ta có  $36\,960$ số.
		\end{itemchoice}
	 }
\end{ex}


%%%%%%%%%%%%%%%% TRẢ LỜI NGẮN %%%%%%%%%%%%%%%%

% nguồn: MaHoaNH_TraLoiNgan_Lop10
\begin{bt}%[0D8H2-3]%[ma:a9faab0e]
	Có bao nhiêu số tự nhiên có ba chữ số dạng $\overline{a b c}$ với $a, b, c \in\{0; 1;  2; 3; 4; 5; 6\}$ sao cho $a<b<c$.
	 \shortans{20}
	\loigiai{ 
		Vì số tự nhiên có ba chữ số dạng $\overline{abc}$ với $a, b, c \in\{0; 1; 2; 3; 4; 5; 6\}$ sao cho $a<b<c$ nên $a,  b, c \in\{1; 2; 3; 4; 5; 6\}$.\\
	Suy  ra số các số có dạng $\overline{abc}$ là $\mathrm{C}_6^3=20$.}
\end{bt}

% nguồn: MaHoaNH_TraLoiNgan_Lop10
\begin{bt}%[0D8H2-3]%[ma:10160104]
	Từ tập $X=\{2; 3; 4; 5; 6\}$ có thể lập được bao nhiêu số tự nhiên có ba chữ số mà các chữ  số đôi một khác nhau?
	 \shortans{60}
	 \loigiai{
		Số các số tự nhiên có ba chữ số mà các chữ số đôi  một khác nhau được lập từ tập $X$ là số chỉnh hợp chập $3$ của $5$ phần tử.\\
	Vậy số các số cần lập  là $\mathrm{A}_5^3=60$ (số).}
\end{bt}

% nguồn: MaHoaNH_TraLoiNgan_Lop10
\begin{bt}%[0D8H2-3]%[ma:508beef9]
	Từ các chữ số $0; 1; 2;  3; 4; 5; 6$ có thể lập được bao nhiêu số tự nhiên có ba chữ số khác nhau?
	 \shortans{180}
	\loigiai{ 
		 Số cách chọn ra chữ số hàng trăm là $6$ cách.\\
		 Với chữ số hàng chục và chữ số hàng đơn vị, mỗi cách chọn ra $2$ số chính là một chỉnh hợp chập $2$ của $6$ phần tử.\\
	Vậy số các số tự nhiên có ba chữ số khác nhau  lập được là $6 \cdot \mathrm{A}_6^2=180$.}
\end{bt}

% nguồn: MaHoaNH_TraLoiNgan_Lop10
\begin{bt}%[0D8H2-3]%[ma:a21f3328]
	Từ các chữ số $1$; $2$; $3$; $4$; $5$; $6$; $7$; $8$. Có thể lập được bao nhiêu số tự nhiên chẵn  có bốn chữ số khác nhau?\shortans{840}
	\loigiai{ 
		Gọi số lập được là  $x = \overline{abcd}$.\\
		Vì  $x$ chẵn nên $d \in \left\lbrace 2;4;6;8\right\rbrace $ nên chọn $d$ có $4$ cách.\\
		Chọn $\overline{abc}$ có $\mathrm{A}_{7}^{3}$  cách.\\
		Vậy có $4\cdot \mathrm{A}_{7}^{3} = 840$ số  thỏa mãn yêu cầu bài toán.
	} 
\end{bt}

% nguồn: MaHoaNH_TraLoiNgan_Lop10
\begin{bt}%[0D8H2-3]%[ma:cfbed44d]
	Từ các số $1$; $2$; $3$; $4$; $5$; $6$; $7$ lập được bao nhiêu số tự nhiên có $4$ chữ số phân biệt  trong đó có $2$ chữ số lẻ và $2$ chữ số chẵn?
	 \shortans{432}
	\loigiai{ 
		 Chọn $2$ chữ số lẻ từ $7$ chữ số đã cho có $\mathrm{C}_{4}^{2}$ cách.\\
		Chọn $2$ chữ số chẵn từ $7$ chữ số  đã cho có $\mathrm{C}_{3}^{2}$ cách.\\
		Với $4$ chữ số đã chọn ta xếp vào $4$ vị trí có  $4$! cách.\\
		Do đó có $\mathrm{C}_{4}^{2}\cdot \mathrm{C}_{3}^{2}\cdot 4!=432$ số. 
	 }
\end{bt}

% nguồn: MaHoaNH_TraLoiNgan_Lop10
\begin{bt}%[0D8H2-3]%[ma:4e3a1b10]
	Từ các chữ số $1$; $2$;  $3$; $4$; $5$; $6$; $7$; $8$. Có thể lập được bao nhiêu số tự nhiên có ba chữ số khác nhau và chia hết cho $3$?\shortans{42}
	 \loigiai{
		Gọi số lập được là  $x = \overline{abc}$.\\
		Vì $x \,\vdots\, 3 \Rightarrow a + b +  c \,\vdots\, 3$, do đó
		$$ (a;b;c) \in \left\lbrace (1;2;3), (1;2;6),  (1;3;5), (1;3;8), (1;4;7), (1;5;6), (1;6;8) \right\rbrace. $$
		Mỗi  bộ $3$ số trên ta lập được $3!$ số thỏa mãn.\\
		Vậy số các số lập  được là $7\cdot 3! = 42$ số.
	} 
\end{bt}

% nguồn: MaHoaNH_TraLoiNgan_Lop10
\begin{bt}%[0D8H2-3]%[ma:85eeb198]
	Có bao nhiêu số tự nhiên chia hết cho $2$  mà mỗi số có ba chữ số khác nhau?
	 \shortans{320}
	\loigiai{ 
		Gọi  số có ba chữ số cần tìm là $\overline{abc}$ $(a \neq 0)$.\\
		Vì số cần tìm  chia hết cho $2$ nên $c$ phải là chữ số chẵn.
		\begin{itemize}
			\item TH 1:  $c=0$\\
			Số cách chọn $a$,  $b$ và sắp vào số là $\mathrm{A}_9^2=72$ (cách).\\
			 Do đó trường hợp này có $72$ số.
			\item TH  2: $c\ne0$.\\
			Chọn $c$ có $4$  cách ($2$,$4$,$6$,$8$).\\
			 Chọn $a$ có $8$ cách ($a\ne 0$, $a\ne c$).\\
			 Chọn $b$ có $8$ cách ($b\ne a$, $b\ne c$).\\
			Vậy trường hợp này có $4\cdot  8\cdot8=256$ số.
		\end{itemize}
	Vậy số các số chia hết cho $2$ mà mỗi số có  ba chữ số khác nhau là $72+256=328$ (số).}
\end{bt}

% nguồn: MaHoaNH_TraLoiNgan_Lop10
\begin{bt}%[0D8V2-3]%[ma:4d4e3097]
	Từ các chữ số $0$, $1$, $2$,  $3$, $4$ có thể lập được bao nhiêu số gồm $7$ chữ số, trong đó chữ số $2$ có mặt $3$ lần, các chữ số còn lại có mặt đúng một lần.
	 \shortans{}
	\loigiai{ 
		Xếp các chữ  số $0$, $1$, $2$, $2$, $2$, $3$, $4$ thành một hàng có $\dfrac{7!}{3!}$ cách xếp.\\
		Xếp các chữ số $0$, $1$, $2$, $2$, $2$, $3$, $4$ thành một hàng sao cho chữ  số $0$ đứng đầu có $\dfrac{6!}{3!}$ cách xếp.\\
		 Vậy có $\dfrac{7!}{3!}-\dfrac{6!}{3!}=1440$ số thỏa mãn yêu cầu.
	 }
\end{bt}

% nguồn: MaHoaNH_TraLoiNgan_Lop10
\begin{bt}%[0D8V2-3]%[ma:87dbf898]
	Từ các chữ số $1$, $2$, $3$,  $4$, $5$, $6$ lập được bao nhiêu số tự nhiên gồm tám chữ số sao cho trong mỗi số đó có đúng ba chữ số $1$, các chữ số còn lại đôi một khác nhau và hai chữ số chẵn không đứng cạnh nhau?
	\shortans{2400} 
	 \loigiai{
		\begin{itemize}
			\item \textbf{Bước 1:} Ta xếp các số lẻ: có  các số lẻ là $1$, $1$, $1$, $3$, $5$ vậy có $\dfrac{5!}{3!}$ cách xếp.
			\item \textbf{Bước 2:}  Ta xếp $3$ số chẵn $2$, $4$, $6$ xen kẽ $5$ số lẻ trên có $6$ vị trí để xếp $3$ số vậy có $\mathrm{A}_{6}^{3}$ cách xếp.
		\end{itemize}
		Vậy có $\dfrac{5!}{3!}\cdot \mathrm{A}_{6}^{3}=2400$  thỏa mãn yêu cầu bài toán.
	} 
\end{bt}

% nguồn: MaHoaNH_TraLoiNgan_Lop10
\begin{bt}%[0D8V2-3]%[ma:916ef479]
	Từ các chữ số  của tập hợp $\{0;1;2;3;4;5\}$, có thể lập được bao nhiêu số tự nhiên có $5$ chữ số đôi một khác nhau mà trong đó nhất thiết phải có mặt chữ số $0$?
	\shortans{480} 
	\loigiai{ 
		Giả sử số tự nhiên có $5$  chữ số đôi một khác nhau có dạng: $\overline{a_{1}a_{2}a_{3}a_{4}a_{5}}$.\\
		Chọn một số cho  $a_{1}$ ta có $5$ cách chọn.\\
		Tiếp theo ta bỏ số $a_{1}$ và số $0$ thì từ tập hợp đã cho chúng ta  còn lại $4$ số. Ta chọn $3$ số từ $4$ số đó ta có $\mathrm{C}_{4}^{3}$ cách chọn.\\
		Chúng ta xếp  số $0$ và $3$ số vừa mới chọn vào $4$ vị trí $a_{2}$, $a_{3}$, $a_{4}$, $a_{5}$ ta được $4!$ cách xếp.\\
		Chọn cho các số cho $a_{2},a_{3},a_{4},a_{5}$ có mặt  chữ số $0$ ta có $\mathrm{C}_{5}^{3}\cdot 4!$ cách chọn.\\
		Số số tự nhiên thỏa yêu cầu đề bài có thể lập được là $5\cdot 4!\cdot\mathrm{C}_{4}^{3}=480$. 
	 }
\end{bt}

% nguồn: MaHoaNH_TraLoiNgan_Lop10
\begin{bt}%[0D8V2-3]%[ma:cba550e5]
	Cho các chữ số $0$; $2$; $3$; $4$;  $5$; $7$; $8$. Từ các chữ số đó có thể lập được bao nhiêu số có $4$ chữ số khác nhau chia hết cho $20$ và luôn xuất hiện chữ số $4$.
	 \shortans{36}
	 \loigiai{
		Vì số có  $4$ chữ số khác nhau chia hết cho $20$ thì hai chữ số cuối cùng phải chia hết cho $20$ nên suy ra $d=0$. Vậy gọi số có $4$ chữ số khác nhau là $\overline{abc0}$.
		\begin{itemize}
			\item TH 1: Nếu $c=4$.\\ 
			chọn $c$ có $1$ cách chọn và chọn $a$, $b$ có $\mathrm{A}_{5}^{2}=20$  cách.
			 \item TH 2: Nếu $c\neq 4$.\\
			chọn $c$ có $2$  cách chọn và đưa số $4$ vào $2$ vị trí $a$, $b$ có $2$ cách.
		\end{itemize}
		Sau khi đưa $4$  vào một trong hai vị trí $a$ hoặc $b$ thì còn $4$ số đưa vào một vị trí còn lại.\\
		Theo quy tắc  nhân $2\cdot 2\cdot 4=16$ cách.\\
		Vậy theo  quy tắc cộng $20+16=36$.
	 }
\end{bt}

% nguồn: MaHoaNH_TraLoiNgan_Lop10
\begin{bt}%[0D8V2-3]%[ma:08e86edc]
	Từ các chữ số $0$; $1$; $2$; $3$; $4$; $5$ có  thể lập được bao nhiêu số tự nhiên gồm $8$ chữ số trong đó chữ số $1$ xuất hiện $3$ lần, các chữ số còn lại xuất hiện đúng $1$ lần.
	\shortans{5880} 
	\loigiai{ 
		Gọi $\overline{n}=\overline{a_{1}a_{2}\ldots a_{8}}$ thỏa  mãn đề bài.
		\begin{itemize}
			\item \textbf{Trường hợp 1} Nếu $a_{1}=1$ thì có $\mathrm{C}_{7}^{2}$ cách sắp $2$ chữ số $1$ vào  $7$ vị trí $\overline{a_{2}\ldots a_{8}}$.\\
			Có  $5$! cách sắp $5$ chữ số $0$; $2$; $3$; $4$; $5$.\\
			Trường hợp này có  $\mathrm{C}_{7}^{2}\cdot 5$! số.
			\item \textbf{Trường hợp 2} Nếu $a_{1}\neq  1$ thì có $\mathrm{C}_{7}^{3}$ cách sắp $3$ chữ số $1$ vào $7$ vị trí $\overline{a_{2}\ldots a_{8}}$.\\
			Có $4$ cách chọn $a_{1}$, có $4$! cách sắp $4$ còn lại.\\ 
			Trường hợp này có $\mathrm{C}_{7}^{3}\cdot  4\cdot 4!$ số.
		\end{itemize}
		 Vậy có $\mathrm{C}_{7}^{2}\cdot 5!+\mathrm{C}_{7}^{3}\cdot 4\cdot 4!=5880$ số.
	 }
\end{bt}

% nguồn: MaHoaNH_TraLoiNgan_Lop10
\begin{bt}%[0D8V2-3]%[ma:a1f9d5b4]
	Cho tập hợp $A=\{0;1;2;3;4;5;6;7\}$. Có bao nhiêu số  tự nhiên chẵn có $6$ chữ số khác nhau được lập thành từ các chữ số của tập $A$, đồng thời phải có mặt ba chữ số $0;1; 2$ và chúng đứng cạnh nhau?
	 \shortans{624}
	\loigiai{ 
		Gọi số cần tìm có  dạng $\overline{a_1 a_2 a_3 a_4 a_5 a_6}$.\\
		 Ta xét các trường hợp
		\begin{enumerate}
			\item Trường hợp 1: $a_6=0$, suy ra $a_6$ có $1$ cách chọn.\\ 
			Xếp các chữ  số $1;2$ vào vị trí $a_4$ và $a_5$ có $2$ cách.\\
			Chọn thứ tự $a_1;  a_2; a_3$ từ tập $\{3;4;5;6;7\}$ có $\mathrm{A}_5^3$ cách.\\
			 Suy ra trường hợp này có $1\cdot 2\cdot \mathrm{A}_5^3=120$ số.
			\item Trường hợp 2: $a_6=2$. Tương tự như trường  hợp 1 nên có $120$ số.
			\item Trường hợp 3: $a_6\in \{4;6\}$, suy  ra $a_6$ có $2$ cách chọn.\\
			Xếp các  chữ số $0;1;2$ đứng cạnh nhau có $3\cdot 3!-2!=16$ cách.\\
			Chọn thứ tự hai chữ số từ tập $\{3;4;5;6;7\}\setminus\{a_6\}$ để xếp vào hai vị trí còn lại có $\mathrm{A}_4^2$ cách. Do đó trường hợp này  có $2\cdot 16\cdot \mathrm{A}_4^2=384$ số.
		\end{enumerate}
		Vậy có $120+120+384=624$ số. 
	 }
\end{bt}

% nguồn: MaHoaNH_TraLoiNgan_Lop10
\begin{bt}%[0D8V2-3]%[ma:b801c1ab]
	Gọi $S$ là tập tất cả các số tự nhiên có $4$ chữ số đôi một khác nhau và không có hai chữ số $2$ và $3$ đứng cạnh nhau, được tạo từ tập $A=\{0;1;2;3;4;5\}$. Số phần  tử của tập $S$ là
	 \shortans{240}
	 \loigiai{
		Gọi $\overline{a_1a_2a_3a_4}$ là số tự nhiên có $4$ chữ  số khác nhau.\\
		$a_1$ có $5$ cách chọn, $\overline{a_2a_3a_4}$ có  $\mathrm A_3^5=60$ cách chọn.\\
		Suy ra có $ 5\cdot  60=300$ số tự nhiên có $4$ chữ số khác nhau. \\
		Cho $2$ và $3$  đứng cạnh nhau nên ta coi như số $m$.\\
		Tập $A$ lúc này là  $A=\{0;1;m;4;5\}$.\\
		Khi đó $m$ có $2$ cách  chọn.\\
		\textbf{Trường hợp 1:} Số tự nhiên $4$ chữ  số có dạng $\overline{mab}$.\\
		 $a$ có $4$ cách chọn.\\
		 $b$ có $3$ cách chọn.\\
		Suy ra  có $4\cdot3=12$ số. \\
		\textbf{Trường hợp 2:} Số tự nhiên $4$ chữ số có dạng  $\overline{amb}$.\\
		$a$ có $3$  cách chọn.\\
		$b$  có $3$ cách chọn.\\
		Suy ra có $3\cdot3=9$ số.  \\
		\textbf{Trường hợp 3:} Số  tự nhiên $4$ chữ số có dạng $\overline{abm}$.\\
		 $a$ có $3$ cách chọn.\\
		$b$ có $3$ cách  chọn.\\
		Do đó có $3\cdot3=9$ số.  \\
		Suy ra có $2(12+9+9)=60$ số tự nhiên có $4$ chữ  số đôi một khác nhau và có hai chữ số $2$ và $3$ đứng cạnh nhau.\\
		Vậy số phần tử của  tập $S$ là $300-60=240$.
	} 
\end{bt}

% nguồn: MaHoaNH_TraLoiNgan_Lop10
\begin{bt}%[0D8C2-3]%[ma:d0f3db3b]
	Có bao nhiêu số tự  nhiên gồm $4$ chữ số đôi một khác nhau, chia hết cho $15$ và mỗi chữ số đều không vượt quá $5$?
	 \shortans{38}
	\loigiai{ 
		Mỗi chữ số đều không vượt quá  $5$. Ta lập số từ tập hợp $\{0;1;2;3;4;5\}$\\
		Số chia hết cho $15$ là  số vừa chia hết cho $3$ vừa chia hết cho $5$. Do đó tận cùng nó là $0$ hoặc $5$.
		\begin{itemize}
			\item \textbf{Trường hợp  1:} Số cần lập có dạng $\overline{abc0}$ với $a;b;c\in\{1;2;3;4;5\}$.\\
			Tổng $a+b+c+0$ phải  chia hết cho $3\Rightarrow a+b+c$ chia hết cho $3$.\\
			Có $4$ tập  hợp $\{a;b;c\}$ có tổng các phần tử chia hết cho $3:\{1;2;3\};\{2;3;4\};\{3;4;5\};\{1;3;5\}$.\\
			Suy ra  có $4\cdot 3!=24$ số
			\item \textbf{Trường hợp 2:} Số cần lập có dạng  $\overline{abc5}$ với $a;b;c\in\{0;1;2;3;4\}$\\
			Tổng $a+b+c+5$ phải chia hết cho $3\Rightarrow a+b+c$ chia  cho $3$ dư $1$.\\
			Có $3$ tập hợp $\{a;b;c\}$ có tổng các phần tử chia $3$  dư $1:\{0;1;3\};\{0;3;4\};\{1;2;4\}$\\
			Có  $2\cdot (3!-2!)+3!=14$ số.
		\end{itemize}
		Vậy có tất cả $24+14=38$ số thỏa đề  bài.
	} 
\end{bt}

% nguồn: MaHoaNH_TraLoiNgan_Lop10
\begin{bt}%[0D8C2-3]%[ma:8d919e85]
	Từ các chữ số $1$, $2$, $3$, $4$, $5$, $6$, $7$, $8$, $9$  có thể lập được bao nhiêu số tự nhiên có $4$ chữ số đôi một khác nhau và không có hai chữ số liên tiếp nào cùng chẵn?
	\shortans{1800} 
	\loigiai{ 
		Trong các chữ số 1, 2,  3, 4, 5, 6, 7, 8, 9 có 4 chữ số chẵn và 5 chữ số lẻ.\\
		Ta xét các trường hợp sau 
		\begin{itemize}
			\item Trường hợp 1. Số gồm 4 chữ số lẻ, có $\mathrm{A}_5^4= 120$ (số). 
			\item Trường hợp 2. Số gồm 1 chữ số chẵn và 3 chữ số  lẻ.\\
			Chọn 1 chữ số chẵn, có  $\mathrm{C}_4^1$ cách.\\
			Chọn 3  chữ số lẻ, có $\mathrm{C}_5^3$ cách.\\
			Sắp xếp 4 chữ số đó  vào 4 vị trí, có $4!$ cách.\\
			Theo quy tắc nhân, có $\mathrm{C}_4^1 \cdot \mathrm{C}_5^3 \cdot 4!=  960$ (số).
			\item Trường hợp 3. Số gồm 2 chữ số chẵn và 2 chữ số lẻ.\\ 
			Chọn 2 chữ số chẵn, có  $\mathrm{C}_4^2$ cách.\\
			Chọn 2 chữ số lẻ, có  $\mathrm{C}_5^2$ cách.\\
			Sắp xếp 2 chữ số lẻ vào 2 vị trí, có  $2!$ cách.\\
			Để không có 2 chữ số liên tiếp nào cùng chẵn, cần sắp xếp 2 chữ số chẵn vào 2 trong 3 vị trí (trước 2 chữ số lẻ, giữa 2 chữ số lẻ, sau  2 chữ số lẻ), có $\mathrm{A}_3^2$ cách.\\
			Theo quy tắc nhân, có $\mathrm{C}_4^2 \cdot  \mathrm{C}_5^2 \cdot 2! \cdot \mathrm{A}_3^2= 720$ số.
		\end{itemize}
		Vậy theo quy  tắc cộng, số các số thoả mãn bài toán là $120+960+720= 1800$ số.
	} 
\end{bt}


%%%%%%%%%%%%%%%% TỰ LUẬN %%%%%%%%%%%%%%%%

% nguồn: MaHoaNH_TuLuan_Lop10
\begin{bt}%[0D8H2-3]%[ma:cfb56854]
	Có bao nhiêu số tự  nhiên có $5$ chữ số phân biệt thuộc tập $\{1;2;3;4;5\}$?
	 \loigiai{
		Mỗi  số tự nhiên có $5$ chữ số phân biệt được lấy từ tập $\{1;2;3;4;5\}$ là một hoán vị của $5$ phần tử.\\
		Vậy  có $\mathrm{P}_5 = 5! = 120$ số.
	} 
\end{bt}

% nguồn: MaHoaNH_TuLuan_Lop10
\begin{bt}%[0D8H2-3]%[ma:d2c1b394]
	Cho tập hợp $S = \{1; 2; 3; 4\}$. Có bao nhiêu số tự nhiên  có $4$ chữ số phân biệt lấy từ tập $S$?
	\loigiai{ 
		Gọi $\overline{a_1a_2a_3a_4}$ là số cần  tìm, $a_i \in S$, $\forall i = \overline{1; 4}$.\\
		Mỗi hoán vị của $4$ phần tử tập hợp $S$ ta  được $1$ số tự nhiên có $4$ chữ số cần tìm, ví dụ như $x = 3214$.\\
		Do vậy, ta được $\mathrm{P}_4 = 4! =  4 \cdot 3 \cdot 2 \cdot 1 = 24$ số.
	 }
\end{bt}

% nguồn: MaHoaNH_TuLuan_Lop10
\begin{bt}%[0D8H2-3]%[ma:b10c0f09]
	Có bao nhiêu số tự nhiên có $5$ chữ số  phân biệt thuộc tập $\{0;1;2;3;4\}$?
	\loigiai{ 
		Số các số tự nhiên có $5$ chữ số phân biệt thuộc tập $\{0;1;2;3;4\}$ (kể cả số $0$ đứng đầu) là: $\mathrm{P}_5 = 5!  = 120$ số.\\
		Số các số tự nhiên có $5$ chữ số phân biệt thuộc  tập $\{0;1;2;3;4\}$ có số $0$ đứng đầu là: $\mathrm{P}_4 = 4! = 24$ số.\\
		Vậy có $120 - 24 = 96$ số thỏa  mãn yêu cầu bài toán.
	} 
\end{bt}

% nguồn: MaHoaNH_TuLuan_Lop10
\begin{bt}%[0D8H2-3]%[ma:6c4a4dbf]
	Có bao nhiêu số tự nhiên gồm $3$ chữ số khác nhau,  biết tổng của $3$ chữ số này bằng $18$?
	\loigiai{ 
		Gọi  số cần tìm là $\overline{abc}$, với $a \neq 0$.\\
		Từ tập $A = \{0; 1; 2; 3; 4; 5; 6; 7; 8; 9\}$, ta có những tập con của $A$ gồm $3$ phần tử sao cho tổng của chúng bằng  $18$ là:
		\[\{9; 8; 1\}, \{9; 6; 3\}, \{9; 5; 4\}, \{8; 7;  3\}, \{8; 6; 4\}, \{7; 6; 5\}, \{2; 7; 9\}.\]
		Vậy có $7$ tập con có $3$ phần tử thuộc $A$ sao cho tổng của $3$ phần  tử này bằng $18$.\\
		Hoán vị $3$ phần tử trong $1$ tập con này ta  được một số cần tìm. Suy ra có tất cả $3! \cdot 7 = 42$ số thỏa yêu cầu.
	} 
\end{bt}

% nguồn: MaHoaNH_TuLuan_Lop10
\begin{bt}%[0D8H2-3]%[ma:4a1e8476]
	Có bao nhiêu số  tự nhiên chẵn có $5$ chữ số đôi một khác nhau sao cho trong mỗi số đó nhất thiết phải có mặt chữ số $0$?
	 \loigiai{
		Gọi số cần tìm  là: $\overline{a_1a_2a_3a_4a_5}$.\\
		 \textbf{Trường hợp 1:} $a_5 = 0$\\
		Khi đó có $1$  cách chọn $a_5$.\\
		Có $\mathrm{A}_9^4$ cách chọn ra $4$ trong $9$ số còn lại và xếp vào các vị trí  còn lại.\\
		Suy ra  có: $1 \cdot \mathrm{A}_9^4 = \mathrm{A}_9^4$ số.\\
		\textbf{Trường hợp 2:} $a_5 \in  \{2; 4; 6; 8\}$\\
		Khi đó có $4$ cách chọn $a_5$.\\ 
		Xếp chữ số $0$ vào $1$ trong $3$ vị trí $a_2$;  $a_3$; $a_4$ có $3$ cách.\\
		 Có $\mathrm{A}_8^3$ cách chọn $3$ số trong $8$ số còn lại và xếp vào các vị trí còn lại.\\
		Suy  ra có: $4 \cdot 3 \cdot \mathrm{A}_8^3$ số.\\
		Vậy có $\mathrm{A}_9^4 + 4 \cdot 3 \cdot \mathrm{A}_8^3 = 7\,056$ số  thỏa mãn yêu cầu bài toán.
	 }
\end{bt}

% nguồn: MaHoaNH_TuLuan_Lop10
\begin{bt}%[0D8H2-3]%[ma:44b9e6c2]
	Từ các chữ số $0$, $1$, $2$, $3$, $4$, $5$. Hỏi có thể  lập được bao nhiêu số tự nhiên gồm bốn chữ số phân biệt và chia hết cho $4$?
	\loigiai{ 
		Gọi số tự nhiên cần tìm $\overline{a  b c d}$.\\
		Nhận xét: Một số tự nhiên (gồm nhiều chữ số) chia hết cho $4$ khi hai chữ số cuối của nó hình thành một số  tự nhiên chia hết cho $4$.\\
		Theo  đề, ta có $\overline{c d}\in\{04,12,20,24,32,40,52\}$.
		\begin{enumerate}[TH1:]
			 \item $\overline{c d}\in\{04,20,40\}$, có $3$ cách chọn $\overline{cd}$.\\
			 Chọn $a$ có $4$ cách; chọn $b 3$ cách.\\
			 Vậy số các số thỏa mãn là $3\cdot4\cdot3=36$ số.
			\item $\overline{c d}\in\{12,24,32,52\}$, có $4$ cách  chọn.\\
			Chọn $a$ có $3$ cách; chọn  $b$ có $3$ cách. Số các số thỏa mãn là $4\cdot3\cdot 3=36$ số.
		\end{enumerate}
		Vậy số các số tự nhiên thỏa đề  bài là $36+36=72$ số.
	} 
\end{bt}

% nguồn: MaHoaNH_TuLuan_Lop10
\begin{bt}%[0D8H2-3]%[ma:1c9f0b99]
	\begin{enumerate}
		\item Có bao nhiêu số tự nhiên có $5$ chữ số đôi một khác nhau? 
		\item Có bao nhiêu số tự nhiên  có $5$ chữ số và số đó là số chẵn?
		\item Có bao nhiêu số tự nhiên có $5$ chữ số đôi một khác nhau và số đó là số lẻ? 
	\end{enumerate}
	\loigiai{ 
		\begin{enumerate}
			\item Gọi $M = \overline{abcde}$, $(a \neq 0)$ là số có $5$ chữ  số khác nhau.\\
			Ta có $a$ có $9$ cách chọn nên có $\mathrm{A}_9^4$ cách chọn $4$ số xếp vào $4$ vị  trí $\overline{bcde}$.\\
			 Vậy có $9\cdot \mathrm{A}_9^4 = 27216$ số.
			\item Gọi $A = \overline{abcde}$ là số có  $5$ chữ số và $A$ là số chẵn.
			Ta có $a$ có $9$ cách chọn; chọn $b$; $c$; $d$ mỗi số có $10$ cách chọn; $e$ có $5$ cách  chọn.\\
			Vậy có $9\cdot  10^3\cdot 5 = 45000$ số.
			\item  Gọi $B = \overline{abcde}$ là số có $5$ chữ số và $B$ là số lẻ.\\
			Ta có $e$ có $5$ cách chọn; $a$ có $8$ cách chọn; có $\mathrm{A}_8^3$ cách chọn chữ số xếp vào ba vị  trí $b$; $c$; $d$.\\
			Vậy có $5\cdot 8\cdot \mathrm{A}_8^3 =  13440$ số.
		\end{enumerate}
	} 
\end{bt}

% nguồn: MaHoaNH_TuLuan_Lop10
\begin{bt}%[0D8V2-3]%[ma:e1fb4b3b]
	Tính số các số tự nhiên đôi một khác nhau có $6$ chữ số  tạo thành từ các chữ số $0,1,2,3,4,5$ sao cho hai chữ số $3$ và $4$ đứng cạnh nhau.
	 \loigiai{
		Xét  số có hình thức $\overline{0bcdef}$.\\
		Số cách hoán đổi vị trí  hai chữ số $3$, $4$ (cùng nhóm $X$) là $2$.\\
		Số cách hoán đổi vị trí của $X$ với các chữ  số $1$, $2$, $5$ là $4!$.\\
		Do đó, số các số được  lập theo hình thức này là $2\cdot 4!=48$.\\
		Xét số có  hình thức $\overline{abcdef}$ trong đó $a$ được phép bằng $0$.\\
		Số cách hoán đổi vị  trí của hai chữ số $3$, $4$ (cùng nhóm $X$) là $2$.\\
		Số cách hoán đổi vị trí của $X$ với các chữ số $0$, $1$,  $2$, $5$ là $5!$\\
		Số các số được lập theo hình thức này là $2\cdot  5!=240$.\\
	Vậy số các  số tự nhiên thỏa mãn đề bài là $240-48=192$.}
\end{bt}

% nguồn: MaHoaNH_TuLuan_Lop10
\begin{bt}%[0D8C2-3]%[ma:770c4428]
	Có bao nhiêu số tự nhiên có $5$ chữ số mà tổng các chữ số trong mỗi  số là $3$.
	 \loigiai{
		 Tổng $5$ chữ số bằng $3$ thì tập hợp các số đó có thể là $\{0; 1; 2\},\{1; 1; 1\},\{3; 0\}$.\\
		{\bf Trường hợp 1.} Số có $5$ chữ số  gồm $3$ chữ số $0$, $1$ chữ số $1$ và $1$ chữ số $2$.\\
		Chọn chữ số xếp vào vị trí đầu có $2$ cách, xếp chữ số còn lại vào $4$ vị trí  cuối có $4$ cách.\\
		Do đó, trường hợp này có $2\cdot 4=8$  (số).\\
		{\bf Trường hợp 2.} Số  có $5$ chữ số gồm $3$ chữ số $1$, $2$ chữ số $0$.\\
		Xếp số $1$ vào vị  trí đầu có $1$ cách.\\
		Xếp $2$ số $1$ còn lại  vào $4$ vị trí cuối có $\mathrm{C}_4^2$ cách.\\
		Do  đó, trường hợp này có $\mathrm{C}_4^2=6$ (số).\\
		{\bf Trường hợp 3.} số có $5$ chữ số gồm $1$  chữ số $3$, $4$ chữ số $0$. Trường hợp này duy nhất $1$ số là $30\ 000$.\\
	Vậy có $8+6+1=15$ số thoả yêu  cầu bài toán.}
\end{bt}
